\section{Ricerca binaria sul filesystem}
\begin{frame}[fragile]{find - ricerca binaria sul filesystem}
  IDEA: Non cercare contenuti, cerca i contenitori.
\begin{figure}
\begin{lstlisting}
find [PERCORSO] [FILTRI] [AZIONE]
\end{lstlisting}
\end{figure}

	\begin{itemize}
		\item PERCORSO: L'origine della ricerca (. per la cartella corrente, / per la root).
		\item FILTRI: I criteri di selezione dei file.
		\item AZIONE: Cosa eseguire sui risultati (se non specificato, li stampa a schermo).
	\end{itemize}
\end{frame}

\begin{frame}[fragile]{find - filtri di ricerca}
 FILTRI: I criteri di selezione dei file.
		\begin{itemize}
			\item Per Tempo (Usa i segni: - meno di, + più di)
			\begin{itemize}
				\item mtime (Giorni dall'ultima modifica del contenuto).
				\item mmin (Minuti dall'ultima modifica).
				\item atime / ctime (Accesso o modifica metadati).
			\end{itemize}
			\item Per Dimensione
			\begin{itemize}
				\item -size (Es. +1G per "più di 1 GB", -50k per "meno di 50 KB").
				\item -empty (Trova file o cartelle completamente vuoti).
			\end{itemize}
			\item Per Utente / Permessi
			\begin{itemize}
			\item -user (Cerca per proprietario del file).
			\item -perm (Cerca permessi esatti come 777 o parziali come /u=w).
			\end{itemize}
		\end{itemize} 	
\end{frame}

\begin{frame}[fragile]{find - operatori logici}
	Esistono anche operatori logici per combinare più filtri e creare ricerche più complesse.
\begin{itemize}
	\item AND (Implicito): Basta scriverli di seguito (es. file .log E più grandi di 10MB).
	\item OR (-o): Ne basta uno vero (es. file .pdf O .doc).
	\item NOT (!): Inverte il match (es. tutto tranne i file .txt).
\end{itemize}
\end{frame}

\begin{frame}[fragile]{find - il vero potere in azione}
	È possibile manipolare i risultati di find usando l'opzione di esecuzione dei comandi sui file trovati (\texttt{-exec})
 \begin{itemize}
	\item \texttt{\{\}} Il segnaposto che rappresenta il file trovato.
	\item \texttt{;} Chiude il comando (lo esegue una volta per ogni file).
	\item \texttt{+} Chiude il comando accodando tutti i file (lo esegue una volta sola in blocco, molto più efficiente).
	\item \texttt{-delete} Eliminazione diretta dei file trovati (da usare con estrema cautela)
\end{itemize}
	Qualche esempio:
\begin{itemize}
	\item Pulizia automatica file vecchi \code{find /tmp -type f -mtime +30 -delete}
	\item Rende eseguibili tutti gli script in massa \code{find . -name "*.sh" -exec chmod +x {} +}
\end{itemize}
\end{frame}

\begin{frame}[fragile]{find - ottimizzazione e sicurezza}
È importante usare find con attenzione, soprattutto quando si eseguono comandi sui risultati, per evitare di causare danni involontari (es. eliminare file importanti). Alcuni consigli:
\begin{itemize}
	\item Utilizzo della Gestione Spazi (\texttt{-print0}): Stampa i risultati separandoli con un carattere nullo, fondamentale quando si passano i dati a \texttt{xargs -0} per non "rompere" i nomi di file con spazi.
	\item Limiti di Ricerca ({\texttt{-maxdepth}}): 	Es. \texttt{-maxdepth 2} impedisce a find di scendere in un loop infinito di sottocartelle, fermandosi al secondo livello.
\end{itemize}
\medspace
Un esempio: \code{find / -type f -size +100M -atime -30 -exec ls -lh {} \;}
\end{frame}

